\subsection{Solution to Problem 29: Irrationality of Logarithm}

\begin{solution}
Assume, for contradiction, that $\log_n(n+1)$ is rational. Then we can write
\[
\log_n(n+1) = \frac{p}{q}
\]
for some positive integers $p$ and $q$. Converting to exponential form,
\[
n^{p/q} = n + 1
\quad\Longrightarrow\quad
n^p = (n+1)^q.
\]
Since $n$ and $n+1$ are consecutive integers, they are coprime: no prime divides both.

Every prime factor of $n^p$ divides $n$, while every prime factor of $(n+1)^q$ divides $n+1$.
Because $n$ and $n+1$ share no common prime factors, the equality $n^p = (n+1)^q$ can hold only if both sides equal $1$.

But $n > 1$, so $n^p \ge 2^p \ge 2 > 1$ and $(n+1)^q \ge 3^q \ge 3 > 1$. This is a contradiction.

Therefore $\log_n(n+1)$ is irrational.
\end{solution}

\begin{takeaways}
\begin{itemize}
    \item Proof by contradiction: Assume the logarithm is rational and write it as $p/q$
    \item Coprime property: Consecutive integers $n$ and $n+1$ have no common prime factors
    \item Unique factorisation: Equal powers force the same prime factors on both sides
\end{itemize}
\end{takeaways}
